% Build:  tectonic paper/greedy-3sumfree-note.tex
\documentclass[10pt,reqno]{amsart}

\usepackage[T1]{fontenc}
\usepackage[utf8]{inputenc}
\usepackage{amsmath,amssymb,amsthm}
\usepackage{booktabs}
\usepackage{array}
\usepackage{microtype}
\usepackage[margin=0.8in]{geometry}
\usepackage[colorlinks=true,linkcolor=black,citecolor=black,urlcolor=blue]{hyperref}

\theoremstyle{plain}
\newtheorem{theorem}{Theorem}
\newtheorem{lemma}[theorem]{Lemma}
\newtheorem{corollary}[theorem]{Corollary}
\newtheorem{proposition}[theorem]{Proposition}
\newtheorem{conjecture}[theorem]{Conjecture}

\theoremstyle{definition}
\newtheorem{definition}[theorem]{Definition}
\newtheorem{observation}[theorem]{Observation}

\theoremstyle{remark}
\newtheorem{remark}[theorem]{Remark}

\numberwithin{theorem}{section}

% compact vertical spacing
\setlength{\abovedisplayskip}{5pt}
\setlength{\belowdisplayskip}{5pt}
\setlength{\abovedisplayshortskip}{3pt}
\setlength{\belowdisplayshortskip}{3pt}
\makeatletter
\def\thm@space@setup{\thm@preskip=5pt \thm@postskip=5pt}
\makeatother

\newcommand{\Z}{\mathbb{Z}}
\newcommand{\N}{\mathbb{N}}
\newcommand{\G}{G}

\begin{document}

\title[Conjecture 17 on greedy 3-sumfree sequences]{The greedy 3-sumfree meta-conjecture:
a proof, sharp on the range $g \ge d$}

\author{Andrew Hingston}
\email{\texttt{hingstonandrew@icloud.com}}

\date{18 August 2026}

\subjclass[2020]{11B75, 03B25, 68Q45}
\keywords{greedy sum-free sequence, Presburger arithmetic, decision procedure, automatic
sequence, Walnut, Peanut, Journal of Integer Sequences}

\begin{abstract}
Bosma, Bruin, Fokkink, Grube, Reuijl and Tromp \cite{Fokkink} conjecture (Conjecture~17,
their ``meta-conjecture'') a closed periodic description of the greedy 3-sumfree sequence
$S_{1,g,g+d}$ for every $d \ge 2$ and $g \ge d+1$, and report that they could not verify it
uniformly in $(g,d)$ because their decision procedure, Walnut, rejects the required term
$k \cdot(5g+2d)$ as a product of two variables. We prove the conjecture, and show it remains
true down to $g = d$: writing $\G(g,d) = \{1,g\} \cup [g+d,2g+d] \cup
\bigcup_{k \ge 1}\bigl(k(5g{+}2d) + [g{+}d{-}2,\,2g{+}d{-}2]\bigr)$ for the conjectured answer
set, we prove $\G(g,d) = S_{1,g,g+d}$ iff $d \ge 2$, $g \ge d$ and $(g,d) \ne (2,2)$ ---
exactly the range on which it can hold. This is strictly larger than the
conjectured hypothesis $g \ge d+1$: it adds the diagonal $g=d$, $d \ge 3$. The proof is by
hand, an interval-sumset argument in the style O.\ Shtrezi used for the $d=1$ companion
conjecture, and independently by machine: the paper's obstruction is removed by a one-line
change of variables that turns the parametrised-modulus statement into ordinary Presburger
arithmetic, decided by the Peanut decision procedure as eight sentences quantified over all
$(g,d)$ at once (6.2\,s total), and reproduced in Walnut~8-dev. Every finite claim was
independently re-verified, by a second implementation sharing no code with the first, over
thousands of $(g,d)$ cells and by a second decision procedure.
\end{abstract}

\maketitle

\section{Background and the conjecture}
\label{sec:bg}

For integers $x<y<z$, the \emph{greedy 3-sumfree sequence} $S_{x,y,z}$ is the increasing
sequence beginning $x,y,z$ in which each later term is the least integer exceeding the
previous term that is \emph{not} a sum of three distinct earlier terms. $S_{1,2,3}$ and
$S_{1,3,4}$ are OEIS \href{https://oeis.org/A026471}{A026471} and
\href{https://oeis.org/A026475}{A026475}. Bosma, Bruin, Fokkink, Grube, Reuijl and Tromp
\cite{Fokkink} study these sequences with the automatic-sequence decision procedure Walnut,
proving closed forms for $S_{1,2,3}$, $S_{1,3,4}$, $S_{1,4,5}$ (their Theorem~15) and, for the
family $S_{1,g,g+1}$, a conjecture (their Conjecture~16) later proved by O.\ Shtrezi
\cite{Shtrezi}. For the two-parameter family $S_{1,g,g+d}$, $d \ge 2$, they report ``firm
computational evidence, from an implementation in Magma'', for the following, quoted
verbatim:

\begin{conjecture}[{\cite[Conjecture~17]{Fokkink}}]
\label{conj17}
Let $d\geq 2$. For every $g\geq d+1$ the greedy 3-sumfree sequence $S_{1,g,g+d}$ is
characterized as follows:
\[
z\in S_{1,g,g+d}\iff z\in\{1,g,2g+d-1,2g+d\}\ \text{or}\ z\geq g+d\ \text{and}\
z\bmod 5g+2d\in\{g+d-2, g+d-1,\ldots, 2g+d-2\}.
\]
In particular, for $d\geq 2$ and every $g\geq d+1$ after the first $g+3$ entries (in a
preperiod) the sequence $S_{1,g,g+d}$ modulo $5g+2d$ is periodic with period $g+1$.
\end{conjecture}

(The arXiv preprint version~1 \cite{Fokkink} prints $z>g+d$ where the published Journal of
Integer Sequences text has $z\ge g+d$; this is a typo in v1, since $z=g+d$ is a seed of the
sequence and so must be admitted; the paper's own worked example uses $z\ge 8 = g+d$. We use
the published $\ge$ form throughout, as does the statement above.) The authors state that
Walnut cannot verify Conjecture~\ref{conj17} uniformly in $(g,d)$: the term
$z=k\cdot(5g+2d)+w$ needs the operator $*$ applied to two variables ($k$ and $g$, or $k$ and
$d$), which their tool rejects outright, though ``whenever a numerical value is substituted
for $g$ \dots\ all is fine''; they verify individual instances this way. As of August~2026 no
treatment of Conjecture~\ref{conj17} for general $(g,d)$ was found in the literature; the only
citing work, Shtrezi's proof of the $d=1$ case \cite{Shtrezi}, does not mention $d\ge2$ or the
meta-conjecture, and the paper's own Remarks~18 says the $d=1$ family is \emph{not} a special
case of Conjecture~\ref{conj17}. The conjecture was open.

\section{The theorem}
\label{sec:thm}

For $d\ge 2$, $g\ge 2$ write $m=5g+2d$ and let $\G(g,d)$ be the right-hand side of
Conjecture~\ref{conj17}, viz.
\[
\G(g,d) \;=\; \{1,g\}\ \cup\ [g+d,\,2g+d]\ \cup\ \bigcup_{k\ge1}\bigl(km+[g+d-2,\,2g+d-2]\bigr).
\]

\begin{theorem}
\label{thm:main}
Let $d\ge 2$ and $g\ge 2$. Then $\G(g,d)=S_{1,g,g+d}$ if and only if $g\ge d$ and
$2g\ge d+3$; equivalently, if and only if $g \ge d$ and $(g,d)\ne(2,2)$.
\end{theorem}

In one sentence: Conjecture~\ref{conj17} is true exactly as stated, for $d\ge2$, $g\ge d+1$,
and Theorem~\ref{thm:main} shows it remains true on the strictly larger range $d\ge2$,
$g\ge d$, $(g,d)\ne(2,2)$ --- which is exactly the range on which the description \emph{can}
be correct, since Theorem~\ref{thm:main} is an ``if and only if''. The extra range is the
diagonal $g=d$, $d\ge3$ (e.g.\ $S_{1,3,6}$, $S_{1,4,8}$, $S_{1,5,10}$, \dots), which
Conjecture~\ref{conj17}'s hypothesis $g\ge d+1$ excludes; amusingly the paper's own worked
verification, following its Example~19, is the instance $(g,d)=(4,4)$, on that diagonal and
outside the conjecture's own hypothesis.

\begin{corollary}
\label{cor:period}
For $d\ge2$, $g\ge d$, $(g,d)\ne(2,2)$, $S_{1,g,g+d}$ begins with the $g+3$ terms
$1,\,g,\,g+d,\,g+d+1,\,\dots,\,2g+d$ and from there repeats, modulo $m=5g+2d$, the block of
$g+1$ residues $g+d-2,\,g+d-1,\,\dots,\,2g+d-2$, one block per period --- the explicit form of
Conjecture~\ref{conj17}'s ``periodic with period $g+1$''.
\end{corollary}

\section{Proof}
\label{sec:proof}

\subsection*{The obstruction, and its removal}
Fixing $g,d$ numerically, $z=k\cdot m+w$ is an ordinary linear (Presburger) term and
individual instances of Conjecture~\ref{conj17} verify at once, as the source paper does. With
$g,d$ free, $k\cdot(5g+2d)$ is a product of variables and is rejected outright. The fix needs
no new mathematics, only a change of variables: never name $z$, but carry every integer as a
$(\text{quotient},\text{remainder})$ pair against $m$.

\begin{lemma}[Carry split]
\label{lem:carry}
Let $m\ge1$ and let $a=k_am+w_a$, $b=k_bm+w_b$, $c=k_cm+w_c$, $z=k_zm+w_z$ with
$0\le w_a,w_b,w_c,w_z<m$. Then
\[
a+b+c=z \iff \exists\, j\in\{0,1,2\}:\ w_a+w_b+w_c = w_z + jm\ \ \text{and}\ \
k_a+k_b+k_c+j = k_z .
\]
\end{lemma}
\begin{proof}
Put $j=k_z-(k_a+k_b+k_c)$. If $a+b+c=z$ then $jm=(w_a+w_b+w_c)-w_z$, and the right side lies
in $[-(m-1),3m-3]$, forcing $-1<j<3$, i.e.\ $j\in\{0,1,2\}$; this $j$ is unique since $jm$
determines it, and the displayed equalities hold. The converse is immediate.
\end{proof}

Since $j$ ranges over the fixed finite set $\{0,1,2\}$, each disjunct's $jm=j(5g+2d)$ is
\emph{linear} in $g,d$: the product of variables never reappears. This is folklore --- the
schoolbook carry argument for sums of a bounded number $r$ of terms, giving $j\in
\{0,\dots,r-1\}$ in general --- and is not claimed as new mathematics; the contribution is
recognising that it removes exactly the blocker \cite{Fokkink} reports. It is, as the
underlying attack document puts it, ``a reformulation, not a prover capability''.

\subsection*{The pen-and-paper route}
Write $B=[g+d,2g+d]$ ($g+1$ integers) and $A_k=km+[g{+}d{-}2,2g{+}d{-}2]$ for $k\ge1$, so
$\G(g,d)=\{1,g\}\cup B\cup\bigcup_{k\ge1}A_k$. Using $2^X=\{x+y:x,y\in X,\,x<y\}$ and
$3^X=\{x+y+z:x,y,z\in X,\,x<y<z\}$ (as in \cite{Shtrezi}), the complement of $\G$ above the
seeds splits into gaps $\mathrm{Gap}_0=[2g{+}d{+}1,\,6g{+}3d{-}3]$ and
$\mathrm{Gap}_k=km+[2g{+}d{-}1,\,6g{+}3d{-}3]$ for $k\ge1$. Four propositions, each an interval
computation, give Theorem~\ref{thm:main} by induction:

\begin{proposition}[Maximality]
\label{propA}
If $d\ge2$, $g\ge d$, $2g\ge d+3$, every $z>g+d$ with $z\notin\G$ is a sum of three distinct
elements of $\G$.
\end{proposition}
\begin{proposition}[3-sumfreeness]
\label{propB}
If $d\ge2$, $g\ge2$, no element of $\G$ is a sum of three distinct elements of $\G$. (Needs
only $g,d\ge2$, not $g\ge d$.)
\end{proposition}
\begin{proposition}[Seeds]
\label{propC}
If $d\ge2$, $g\ge2$, then $\G\cap[1,g+d]=\{1,g,g+d\}$.
\end{proposition}
\begin{proposition}[Sharpness]
\label{propD}
If $2\le g<d$, then $z=3g+d+2$ satisfies $z>g+d$, $z\notin\G$, and $z$ is not a sum of three
distinct elements of $\G$; hence $\G\ne S_{1,g,g+d}$. The same holds at $(g,d)=(2,2)$, at
$z=14$.
\end{proposition}

Proposition~\ref{propB} is proved by a twelve-case check: $\G$ splits into the four parts
$\{1\},\{g\},B,A_k$, and every sum of three (not necessarily distinct-part) summands is shown
to land inside a $\mathrm{Gap}$, using $2^B=[2g{+}2d{+}1,4g{+}2d{-}1]$,
$3^B=[3g{+}3d{+}3,6g{+}3d{-}3]$ and the analogous identities for the $A_k$; the only place the
hypothesis $d\ge2$ is used is the lower end of the $1+A_i+A_j$ case, which is exactly why
$d=1$ needs a different modulus (Shtrezi's). Proposition~\ref{propA} covers each
$\mathrm{Gap}_0$/$\mathrm{Gap}_k$ by a chain of four sumset intervals whose junctions are
exactly $g\ge d$ and $2g\ge d+3$. Proposition~\ref{propD}'s witness $z=3g+d+2$ falls in
$\mathrm{Gap}_0$ but below every covering interval that needs $g<d$; the step from ``$z$ is
not a 3-sum'' to ``the greedy rule admits it, so $\G\ne S$'' additionally needs $\G$ and $S$
to agree on $[1,z)$, which they do: the only part of $\mathrm{Gap}_0$ below $z$ is
$[2g{+}d{+}1,3g{+}d{+}1]=1+g+B$, covered with \emph{no} hypothesis at all, so the induction
below runs unobstructed up to $z$ for every $g,d\ge2$. Full case-by-case verification of all
twelve rows and both covering chains is in \cite[\S4]{Attack4}.

The four propositions combine by strong induction: Proposition~\ref{propC} settles $[1,g+d]$;
for $z>g+d$, every representation $z=a+b+c$ with $a,b,c\in\G$ has $a,b,c<z$, so by the
induction hypothesis they are exactly the earlier terms of $S_{1,g,g+d}$; then
Proposition~\ref{propA} (if $z\notin\G$) or Proposition~\ref{propB} (if $z\in\G$) tells the
greedy rule what to do, matching $\G$. This proves Theorem~\ref{thm:main}'s ``if''; the ``only
if'' is Proposition~\ref{propD} together with the trichotomy $g<d$, $(g,d)=(2,2)$, or
$g\ge d\wedge 2g\ge d+3$, which is exhaustive for $d,g\ge2$.

\subsection*{The machine route}
The same statement is also settled by decision procedure, with no hand argument. Encode
$k\cdot m+w \in \G(g,d)$ as a predicate $\mathrm{ING}(k,w)$ in $(k,w,g,d)$ directly from the
definition of $\G$; every disjunct forces $w\le 2g+d<m$, so the encoding is injective (checked
as obligation T1 below), which is what licenses Lemma~\ref{lem:carry} without an explicit
$w<m$ guard on the summands. The sum and distinctness conditions become, via
Lemma~\ref{lem:carry}, ordinary linear Presburger formulas in ten free variables
$(g,d,k_a,k_b,k_c,k_z,w_a,w_b,w_c,w_z)$. Propositions~\ref{propA}--\ref{propD} (with a
restricted, sound existential witness family for Proposition~\ref{propA}'s ``there exist
three distinct elements summing to $z$'', since the fully unrestricted six-variable
existential exhausts a 6\,GB budget in both provers after $\sim$280\,s and remains out of
reach) become eight closed first-order sentences T0--T7, each discharged by automaton
compilation in the Peanut decision procedure in 6.2\,s total (peak 1.3\,GB), and independently
reproduced --- posed as universally quantified sentences, which forces the full ten-fold
existential projection Peanut's `witness'-on-open-formula route does not need --- in Walnut
8-dev (52.7\,s total). The full transcript, formulas, and driver are
\texttt{explore/attack4\_engine.py}, \texttt{explore/attack4\_walnut.py},
\texttt{results/attack4\_transcript.txt}, \texttt{results/attack4\_engine.json} and
\texttt{results/attack4\_walnut\_uniform.json}, all in the repository below; the full account
is \cite{Attack4}, \S5.

\section{Verification}
\label{sec:verify}

Every claim above was re-derived and independently re-checked, with code sharing nothing with
the implementation that produced it, in a referee pass whose transcript is \cite{Verdict}. The
hand proof of \S\ref{sec:proof} was re-derived step by step, including all twenty-four
interval-endpoint inequalities and both covering chains, and no gap was found; the primary
source \cite{Fokkink} was independently re-fetched from the published J.\ Integer Seq.\ LaTeX
source, not quoted from memory, and Conjecture~\ref{conj17}'s wording and the $\ge$-vs-$>$
point confirmed against it directly.

\begin{center}\footnotesize
\begin{tabular}{@{}>{\raggedright\arraybackslash}p{6.3cm}>{\raggedright\arraybackslash}p{5.6cm}>{\raggedright\arraybackslash}p{3.3cm}@{}}
\toprule
Check (independent implementation) & Scope & Result \\
\midrule
Theorem~\ref{thm:main}'s iff vs.\ real greedy sequences
  & $d=2\dots30$, $g=2\dots200$, 40 periods of $5g+2d$ (5771 cells)
  & 0 mismatches \\
Proposition~\ref{propD}'s first-failure point $3g+d+2$ (and $14$ at $(2,2)$)
  & every failing cell above
  & exact in all of them \\
adversarial spot checks, e.g.\ $(100,100)$, $(200,3)$, $(3,200)$, $(52,50)$
  & 11 cells
  & 0 mismatches \\
Propositions~\ref{propA}--\ref{propD}, Lemma~\ref{lem:carry}, encoding injectivity, all
  twelve interval containments, both covering chains, restricted-witness soundness and
  completeness
  & $d=2\dots14$, $g=2\dots30$ (377 cells)
  & 0 failures \\
Proposition~\ref{propB} in the failing region $g<d$
  & same grid
  & never fails \\
$d=1$ control, meta-conjecture's own modulus $5g+2$
  & $g=2\dots11$
  & description \emph{and} 3-sumfreeness both fail --- confirms $d\ge2$ is load-bearing \\
\bottomrule
\end{tabular}
\end{center}

The eight obligations of \S\ref{sec:proof} were also re-posed independently, as closed
universally quantified sentences in a different (equivalent) predicate shape, and re-run in
both Peanut and Walnut~8-dev: all TRUE, including the ten-variable 3-sumfreeness sentence
(Peanut, 3.50\,s, peak 1252\,MB; Walnut, 41.8\,s) and, transcribed verbatim from the published
LaTeX source, the paper's own worked $(g,d)=(4,4)$ Walnut fragment (TRUE, 0.3\,s) --- two
independent decision-procedure implementations, two independent formula shapes, agreeing on
every check. The unrestricted six-variable existential for Proposition~\ref{propA} was
attempted again independently and killed by a memory watchdog at 4874\,MB after 274.5\,s,
confirming it is genuinely out of reach with both provers as currently implemented, not an
artefact of one encoding.

\section{Remarks}
\label{sec:remarks}

\subsection*{What is new, and what is not}
\S\ref{sec:proof}'s pen-and-paper argument is a competent transposition of Shtrezi's
interval-sumset technique \cite{Shtrezi} to a different modulus and range, not a new
technique; the partition of $\G$, the gap structure, and the two-proposition induction are
his. What is new is doing this for $d\ge2$, which nobody had done, and the \emph{exact} range:
Proposition~\ref{propD}, which has no counterpart in \cite{Shtrezi}, is what turns an
implication into the sharp equivalence of Theorem~\ref{thm:main} and extends
Conjecture~\ref{conj17} to the diagonal $g=d$, $d\ge3$. Lemma~\ref{lem:carry} is schoolbook
mathematics; its role here is a change of variables, not a theorem, and it should be credited
as such --- the honest place to put the credit is the reformulation, not any capability of the
decision procedure. The explicit preperiod/period of Corollary~\ref{cor:period}, which
\cite{Fokkink} states only as a count, is a small genuine addition.

\subsection*{What is not settled}
The \emph{irregular} region $2\le g<d$ and $(g,d)=(2,2)$ is left open: Theorem~\ref{thm:main}
shows the regular description $\G(g,d)$ fails there, but says nothing about the correct
description, which \cite[Remarks~18]{Fokkink} already notes is irregular (their Example~19,
$S_{1,5,12}$, is periodic with modulus $321$ and period $32$ --- not of the form $5g+2d$).
Start values other than $1$ are untouched. Whether a better decision-procedure ladder can
project the fully unrestricted six-variable existential of Proposition~\ref{propA} directly,
rather than via the restricted-but-sound witness family used here, remains open; \S4
confirms it is out of reach in both Peanut and Walnut~8-dev as they currently stand.

\subsection*{Acknowledgement and declaration of AI usage}
This note, and the work behind it, were produced with substantial use of a large language
model. This note carries this declaration from its first version. Precisely: the models were Claude (Anthropic): \texttt{claude-fable-5} as the orchestrating model, with \texttt{claude-opus-5} and \texttt{claude-sonnet-5} sub-agents for search, implementation, refereeing and drafting, all used through the Claude Code agent framework; they were used for (i) the literature and primary-source search, including
fetching and checking the published Journal of Integer Sequences LaTeX source directly rather
than quoting it from memory; (ii) the design and implementation of Peanut, the open-source
decision procedure with which every machine computation here was run; (iii) the discovery of
the carry-split reformulation of Lemma~\ref{lem:carry}, and the statement and proof of
Theorem~\ref{thm:main} and Propositions~\ref{propA}--\ref{propD}, following the interval-sumset
technique of \cite{Shtrezi}; (iv) an adversarial ``referee'' pass, also carried out by the
model, which re-derived every step by hand, re-read the primary source independently, wrote
brute-force and decision-procedure checks sharing no code with the original implementation,
and found and corrected several defects in an earlier draft of the underlying analysis; and
(v) the drafting of this text, most of which is the model's output lightly edited by the
author. The author directed the work, chose the problem, read and takes responsibility for
every statement, and ran the computations. All finite claims were verified by the independent
implementations described in Section~\ref{sec:verify}; the referee transcript is
\texttt{research/referee-verdicts/attack4-verdict.md} in the repository below. Any errors are the author's.

\begin{thebibliography}{9}
\setlength{\itemsep}{0pt}\setlength{\parsep}{0pt}\setlength{\parskip}{0pt}

\bibitem{Attack4}
A. Hingston,
\emph{Attack target 4 --- the greedy 3-sumfree meta-conjecture},
research note, \texttt{research/attacks/ATTACK-4.md} in the Peanut repository,
\url{https://github.com/whehwhey/peanut}, 2026.

\bibitem{Fokkink}
W. Bosma, R. Bruin, R. Fokkink, J. Grube, A. Reuijl, and T. Tromp,
\emph{Using Walnut to Solve Problems from the OEIS},
J. Integer Seq. \textbf{28} (2025), Article 25.3.8;
arXiv:2503.04122v1 [math.NT], 6 March 2025.

\bibitem{OEIS-A026471}
OEIS Foundation Inc.,
\emph{The On-Line Encyclopedia of Integer Sequences}, entry A026471,
\url{https://oeis.org/A026471}.

\bibitem{OEIS-A026475}
OEIS Foundation Inc.,
\emph{The On-Line Encyclopedia of Integer Sequences}, entry A026475,
\url{https://oeis.org/A026475}.

\bibitem{Shtrezi}
O. Shtrezi,
\emph{The greedy 3-sumfree sequence $S_{1,g,g+1}$},
arXiv:2606.17447 [math.NT], 13 June 2026.

\bibitem{Verdict}
A. Hingston,
\emph{Referee verdict on \texttt{research/attacks/ATTACK-4.md} in the Peanut repository},
\texttt{research/referee-verdicts/attack4-verdict.md},
\url{https://github.com/whehwhey/peanut}, 2026.

\end{thebibliography}

\end{document}
